\chapter{Static Games}\label{ch:b2:static-games}

Four people each receive 20 tokens. Each can keep tokens or put them into a common project; every token in
the project pays 0.4 tokens to each of the four, including the one who put it in. A token kept is worth 1
to its owner, a token contributed only 0.4, so whatever the others do, each person does best by keeping
everything. Yet if all four contribute everything, each ends with 32 tokens instead of 20. The chart shows
what people did when they played this game for money, ten times in a row.

\begin{openingchart}{Mean contribution to the common project, in per cent of the endowment of 20 tokens,
  over ten periods and in the final period: groups re-formed at random each period (Strangers) or kept
  fixed (Partners), without and with the possibility of punishing other members at a cost.}{Fehr and
  G\"achter (2000), Tables~3 and~4}
\begin{tikzpicture}
\begin{axis}[omrecord, ybar, bar width=12pt, xmin=0.5, xmax=4.5, ymin=0, ymax=100, xtick=data,
    xticklabels from table={figdata/bachelor-2/10-static-games/out/contrib.csv}{label},
    xticklabel style={font=\footnotesize, text width=22mm, align=center},
    ylabel={Per cent of the endowment}, grid=none, ymajorgrids, width=0.75\linewidth,
    height=0.45\linewidth, legend columns=2,
    legend style={at={(0.5,-0.3)}, anchor=north, font=\footnotesize}]
  \addplot[fill=omDef!70, draw=omDef] table[x=pos, y=all]
    {figdata/bachelor-2/10-static-games/out/contrib.csv};
  \addplot[fill=omThm!60, draw=omThm] table[x=pos, y=final]
    {figdata/bachelor-2/10-static-games/out/contrib.csv};
  \legend{Mean over ten periods, Final period}
\end{axis}
\end{tikzpicture}
\end{openingchart}

% ledger: b2:fg2000-design, b2:fg2000-t3, b2:fg2000-t4, b2:fg2000-text
% figdata: contrib_summary (strangers 18.5 / 9.5, partners 37.5 / 16, largest 91)
Without punishment, groups of strangers contributed 18.5 per cent of their tokens on average and 9.5 per
cent in the last period; fixed groups contributed 37.5 per cent, falling to 16 per cent. Between 53 and 75
per cent of the players contributed nothing at all in the last period. When players could pay to reduce
another member's earnings after seeing the contributions, the pattern changed completely: fixed groups
contributed 85 per cent on average and 91 per cent in the final period.

This chapter builds the tools to say what a game predicts. A \emph{game} is a situation in which each
person's payoff depends on what others do; the games of this chapter are \emph{static}: everyone chooses
once, at the same time, without seeing the others' choices. The central prediction, \omterm{def:b2:static-games:nash-equilibrium}{Nash equilibrium},
says that no player gains by changing her choice alone. For the game above it predicts zero
contributions, which the first bars approach and the last bars refute. Both results have something to
teach: the first about how incentives shape choices, the second about what the payoffs written into the
game leave out.

\begin{recall}
From the school volumes: the prisoner's dilemma, public goods and the commons. From the Year 1 volume:
Cournot and Bertrand competition and the best-response functions of oligopoly. From
\cref{ch:b2:social-choice}: the median voter.
\end{recall}

\begin{notation}[Symbols of this chapter]\label{not:b2:static-games:symbols}
Players are $i = 1, \dots, n$; $S_i$ is player $i$'s set of strategies and $s_i$ one of them; $s = (s_1, \dots,
s_n)$ is a \omterm{def:b2:static-games:profile}{strategy profile} and $s_{-i}$ the strategies of everyone but $i$; $u_i(s)$ is $i$'s payoff. In the
applications, $q_i$ is a firm's quantity, $g_i$ a contribution with total $G$, $y$ an \omterm{def:b2:exchange-economies:endowment}{endowment}, $e_i$ a
fishing effort with total $E$, and $f$ a fee.
\end{notation}

\section{The normal form}\label{sec:b2:static-games:normal-form}

\begin{definition}[Normal-form game]\label{def:b2:static-games:normal-form}
A \emph{normal-form game}\index{normal-form game} consists of a set of players $i = 1, \dots, n$, a set of
strategies $S_i$ for each player, and for each player a payoff function $u_i(s_1, \dots, s_n)$ that depends
on the strategies of all players.
\end{definition}

\begin{definition}[Strategy profile]\label{def:b2:static-games:profile}
A \emph{strategy profile}\index{strategy profile} is a list $s = (s_1, \dots, s_n)$ with one strategy for
each player. We write $s = (s_i, s_{-i})$ to separate player $i$'s strategy from the others'.
\end{definition}

With two players and few strategies, the game fits in a table: rows are the first player's strategies,
columns the second's, and each cell holds the two payoffs, the row player's first. The figure shows three
games that recur through the chapter.

\begin{omfigure}
\begin{tikzpicture}[econ, font=\footnotesize]
  \foreach \ox/\title/\sa/\sb/\ta/\tb/\p/\q/\r/\s in {
      0/Prisoner's dilemma/C/D/C/D/{3, 3}/{0, 4}/{4, 0}/{1, 1},
      5.2/Coordination/A/B/A/B/{2, 2}/{0, 0}/{0, 0}/{1, 1},
      10.4/Matching pennies/H/T/H/T/{1, -1}/{-1, 1}/{-1, 1}/{1, -1}} {
    \node[font=\small\bfseries] at (\ox+1.3,1.9) {\title};
    \node[payoff] at (\ox+0.8,0.4) {\p};
    \node[payoff] at (\ox+2.4,0.4) {\q};
    \node[payoff] at (\ox+0.8,-0.4) {\r};
    \node[payoff] at (\ox+2.4,-0.4) {\s};
    \node at (\ox-0.25,0.4) {\sa};
    \node at (\ox-0.25,-0.4) {\sb};
    \node at (\ox+0.8,1.1) {\ta};
    \node at (\ox+2.4,1.1) {\tb};
  }
\end{tikzpicture}

{\small Three games in normal form. In each cell, the first number is the row player's payoff, the second
the column player's. In the prisoner's dilemma, D (defect) is better than C (cooperate) for each player
whatever the other does, yet both prefer (C, C) to (D, D). In the \omterm{def:b2:static-games:coordination-game}{coordination game}, each wants to match the
other. In matching pennies, the row player wants to match and the column player to mismatch.}
\end{omfigure}

\section{Dominance and iterated elimination}\label{sec:b2:static-games:dominance}

\begin{definition}[Strictly dominant strategy]\label{def:b2:static-games:dominant}
A strategy $s_i$ is \emph{strictly dominant}\index{strictly dominant strategy} for player $i$ if it gives
her a higher payoff than any other strategy of hers, whatever the other players do: $u_i(s_i, s_{-i}) >
u_i(s'_i, s_{-i})$ for every $s'_i \ne s_i$ and every $s_{-i}$.
\end{definition}

\begin{definition}[Strictly dominated strategy]\label{def:b2:static-games:dominated}
A strategy $s'_i$ is \emph{strictly dominated}\index{strictly dominated strategy} if some other strategy
$s_i$ gives player $i$ a higher payoff whatever the others do: $u_i(s_i, s_{-i}) > u_i(s'_i, s_{-i})$ for
every $s_{-i}$.
\end{definition}

A player who maximises her payoff never plays a \omterm{def:b2:static-games:dominated}{strictly dominated strategy}, whatever she believes the
others will do. In the prisoner's dilemma, C is \omterm{def:b2:static-games:dominated}{strictly dominated} by D for each player, so the prediction
is (D, D), although both would be better off at (C, C). The contribution game of the opening chart is a
prisoner's dilemma with four players and 21 strategies each: since a token contributed returns 0.4 to its
owner, contributing nothing strictly dominates every other choice.

If players also know that the others never play dominated strategies, they can go further.

\begin{definition}[Iterated elimination of strictly dominated strategies]\label{def:b2:static-games:iesds}
\emph{Iterated elimination of strictly dominated strategies}\index{iterated elimination of strictly
dominated strategies} removes every \omterm{def:b2:static-games:dominated}{strictly dominated strategy} of every player, then every strategy that
is \omterm{def:b2:static-games:dominated}{strictly dominated} in the game that remains, and so on until nothing more can be removed.
\end{definition}

The listing applies the procedure to a game with three strategies each, then searches the full game for
profiles in which each player's strategy is a best reply to the other's, the equilibria of the next
section.

\omcode{code/bachelor-2/10-static-games/listings/dominance.py}{1}{31}{\omterm{def:b2:static-games:iesds}{Iterated elimination of strictly
dominated strategies} in a $3 \times 3$ game.}\label{lst:b2:static-games:dominance}

\omoutput{code/bachelor-2/10-static-games/outputs/dominance.txt}{The game is solved by iterated dominance;
the only survivor is also the only pure \omterm{def:b2:static-games:nash-equilibrium}{Nash equilibrium}.}\label{out:b2:static-games:dominance}

% code: listings/dominance.py
Column $C$ goes first: $R$ gives the column player more against every row. Once $C$ is gone, $U$ beats $M$ and
$D$ for the row player; once only $U$ is left, the column player prefers $L$. Each round needs one more level
of ``she knows that he knows'': three rounds here.

\section{Best responses and Nash equilibrium}\label{sec:b2:static-games:nash}

Most games have no dominant strategies: what is best for one player depends on what the others do.

\begin{definition}[Best response]\label{def:b2:static-games:best-response}
A strategy $s_i$ is a \emph{best response}\index{best response} of player $i$ to the others' strategies
$s_{-i}$ if $u_i(s_i, s_{-i}) \ge u_i(s'_i, s_{-i})$ for every $s'_i \in S_i$.
\end{definition}

\begin{definition}[Nash equilibrium]\label{def:b2:static-games:nash-equilibrium}
A \omterm{def:b2:static-games:profile}{strategy profile} $s^* = (s^*_1, \dots, s^*_n)$ is a \emph{Nash equilibrium}\index{Nash equilibrium} if
every player's strategy is a \omterm{def:b2:static-games:best-response}{best response} to the others': $u_i(s^*_i, s^*_{-i}) \ge u_i(s_i, s^*_{-i})$
for every player $i$ and every $s_i \in S_i$.
\end{definition}

In a \omterm{def:b2:static-games:nash-equilibrium}{Nash equilibrium} no player gains by deviating alone. The definition says nothing about how players
reach it: it describes a situation where every player's expectation of the others is correct and she acts
on it. In the prisoner's dilemma the only equilibrium is (D, D). Any profile that survives iterated
dominance alone is a \omterm{def:b2:static-games:nash-equilibrium}{Nash equilibrium} (\cref{exo:b2:static-games:10} asks why the equilibria are never
eliminated).

\begin{definition}[Coordination game]\label{def:b2:static-games:coordination-game}
A \emph{coordination game}\index{coordination game} is a game with several Nash equilibria in which every
player wants to choose the same action as the others, so that players gain from matching their choices.
\end{definition}

The \omterm{def:b2:static-games:coordination-game}{coordination game} of the figure has two equilibria, (A, A) and (B, B), and both players prefer the
first; nothing in the definition says which they will reach. Matching pennies has none in pure
strategies: at every cell one player wants to switch. The existence of a \omterm{def:b2:static-games:nash-equilibrium}{Nash equilibrium} is guaranteed
only once players may randomise, the subject of \cref{ch:b2:mixed-strategies}.

\section{Quantity competition as a game}\label{sec:b2:static-games:cournot}

The Cournot model of the Year 1 volume is a game: the players are firms, the strategies quantities, the
payoffs profits. With inverse demand $p = A - Q$, $Q = \sum_i q_i$, and a constant marginal cost $c$, firm $i$
earns $(A - Q - c)q_i$, and its \omterm{def:b2:static-games:best-response}{best response} to the others' total $Q_{-i}$ is $q_i = (A - c - Q_{-i})/2$.

\begin{proposition}[Cournot equilibrium with $n$ firms]\label{prop:b2:static-games:cournot}
The game has a unique \omterm{def:b2:static-games:nash-equilibrium}{Nash equilibrium}, in which each firm produces $q^* = (A - c)/(n + 1)$ and the price is
$p^* = (A + nc)/(n + 1)$. As $n$ grows, the price falls towards $c$.
\end{proposition}

\begin{proof}
Every firm's quantity must be a \omterm{def:b2:static-games:best-response}{best response}: $2q_i = A - c - Q_{-i}$, that is, $q_i = A - c - Q$ for every
$i$. All firms then produce the same quantity $q$, and $q = A - c - nq$ gives $q^* = (A - c)/(n + 1)$. The price is
$A - nq^* = (A + nc)/(n + 1)$, which tends to $c$ as $n \to \infty$.
\end{proof}

% code: exercises/solutions.py cournot
With $A = 100$ and $c = 10$, a monopoly sells 45 units at 55; two firms sell 30 each at 40; four firms 18 each at 28;
nine firms 9 each at 19. Each firm ignores the fact that its output lowers the price received by the others,
so together they produce more than a monopoly would, and the price falls with their number.

\section{Contributions to a public good}\label{sec:b2:static-games:public-good}

In the opening game the dominant strategy made the prediction stark. With diminishing returns the logic is
the same but the equilibrium interior. Each of $n$ people has income $y = 12$, gives $g_i$ to a public good
$G = \sum_j g_j$ and consumes the rest; her utility is $\ln(y - g_i) + \ln G$. Her \omterm{def:b2:static-games:best-response}{best response} to the others'
total $G_{-i}$ sets $1/(y - g_i) = 1/(g_i + G_{-i})$:
\begin{equation}\label{eq:b2:static-games:br}
g_i = \tfrac12\,(y - G_{-i}), \quad \text{or zero if this is negative.}
\end{equation}

\begin{proposition}[Underprovision]\label{prop:b2:static-games:underprovision}
In this game the \omterm{def:b2:static-games:nash-equilibrium}{Nash equilibrium} is symmetric, $g^* = y/(n + 1)$, and the total $G^* = ny/(n + 1)$ is below the
efficient symmetric total $ny/2$ whenever $n \ge 2$. The ratio of the two, $2/(n + 1)$, falls as the group grows.
\end{proposition}

\begin{proof}
In equilibrium every $g_i$ satisfies \eqref{eq:b2:static-games:br}, so $2g_i + G_{-i} = y$, that is, $g_i + G = y$:
all contributions are equal, and $g + ng = y$ gives $g^* = y/(n + 1)$. Efficiency requires the sum of the marginal
rates of substitution between the public good and private consumption to equal 1, the cost of a unit of the
public good: $\sum_i (y - g_i)/G = 1$. With equal contributions, $n(y - g) = ng$, so $g = y/2$ and the total is
$ny/2$. The ratio is $(ny/(n + 1))/(ny/2) = 2/(n + 1) < 1$ for $n \ge 2$.
\end{proof}

\begin{omfigure}
\begin{tikzpicture}[econ]
\begin{axis}[omrecord, grid=none, xmin=0, xmax=13, ymin=0, ymax=13, xlabel={$g_1$}, ylabel={$g_2$},
    xtick={0,4,6,12}, ytick={0,4,6,12}, width=0.5\linewidth, height=0.5\linewidth]
  \addplot[omDef, very thick, domain=0:12, samples=2] {(12 - x)/2};
  \addplot[omThm, very thick, domain=0:6, samples=2] {12 - 2*x};
  \node[eqpoint] at (axis cs:4,4) {};
  \node[anchor=north east, font=\footnotesize] at (axis cs:3.9,3.9) {Nash};
  \node[eqpoint] at (axis cs:6,6) {};
  \node[anchor=south west, font=\footnotesize] at (axis cs:6.1,6.1) {efficient};
  \node[anchor=south, font=\footnotesize, omDef, align=center] at (axis cs:10.5,2.2) {player 2's\\best response};
  \node[anchor=west, font=\footnotesize, omThm] at (axis cs:0.4,11.5) {player 1's best response};
\end{axis}
\end{tikzpicture}

{\small \omterm{def:b2:static-games:best-response}{Best responses} in the public-good game with two players and $y = 12$. Each line shows one player's best
contribution given the other's. They cross at the \omterm{def:b2:static-games:nash-equilibrium}{Nash equilibrium} (4, 4), with a total of 8; the efficient
symmetric point (6, 6) has a total of 12. Each player gives less than efficiency requires because she counts
only her own benefit from the good.}
\end{omfigure}

% code: exercises/solutions.py public_good
With two players, each gives 4 and the total is 8 instead of 12; with ten, each gives 1.09 and the total of 10.9 is
18 per cent of the efficient 60. The prediction for the linear game of the opening chart is the extreme case:
zero.

\begin{record}[Contributions fall with repetition]
% ledger: b2:fg2000-design, b2:fg2000-t3, b2:fg2000-t4, b2:fg2000-text
Fehr and G\"achter (2000) ran the opening game for ten periods with 24 participants per session. Without
punishment, contributions fell over the periods in both treatments, to 9.5 per cent of the \omterm{def:b2:exchange-economies:endowment}{endowment} among
strangers and 16 per cent in fixed groups by the last period, with between 53 and 75 per cent of
participants contributing nothing. With the option to punish, average contribution rates between 50 and
95 per cent of the \omterm{def:b2:exchange-economies:endowment}{endowment} were maintained, and in the last period of the fixed-group treatment 82.5 per
cent of participants contributed everything.
\end{record}

Punishment costs the punisher and benefits the group, so it is itself a public good: a selfish player
never punishes in the last period, and, knowing this, no one should fear punishment. That the mere option of
punishing raised contributions so much suggests that players expected some others to punish, which a player
who cares only about her own tokens would not do; why this is so, and how cooperation can be
sustained among selfish players who meet again, are questions for \cref{ch:b2:repeated-games}.

\begin{model}{Nash equilibrium as a model of play}
\assumes Players maximise known payoffs, and each correctly anticipates the others' strategies.
\predicts Dominated strategies are not played; zero contributions in a linear public-goods game;
underprovision of public goods and over-use of common resources, worse in larger groups.
\tested Public-goods experiments: contributions start well above zero and fall with repetition, toward
but not to the prediction (Fehr and G\"achter 2000); field studies of commons (Ostrom 1990).
\breaks Players who care about others' payoffs or about fairness; punishment; groups that write and enforce
their own rules; several equilibria with no way to choose among them.
\end{model}

\section{A common resource}\label{sec:b2:static-games:commons}

\begin{definition}[Common-pool resource]\label{def:b2:static-games:common-pool}
A \emph{common-pool resource}\index{common-pool resource} is a resource from which it is costly to exclude
users, and whose use by one person reduces what is left for others: a fishing ground, a pasture, an aquifer.
\end{definition}

A common resource is the mirror image of a public good: each user imposes a cost on the others instead of
conferring a benefit. Let $n$ boats fish one ground. A boat's effort $e_i$ (boat-days) lands a catch worth $A - E$
per boat-day, where $E = \sum_j e_j$: the more the ground is fished, the less each day yields. A boat-day costs $c$.
Boat $i$ earns $e_i(A - E - c)$, and its \omterm{def:b2:static-games:best-response}{best response} sets $A - E - e_i - c = 0$. This is the Cournot game again:
each boat ignores the fall in the others' yield that its effort causes.

\begin{proposition}[Over-use of a common resource]\label{prop:b2:static-games:commons}
In the symmetric \omterm{def:b2:static-games:nash-equilibrium}{Nash equilibrium}, total effort is $E^{\mathrm{N}} = n(A - c)/(n + 1)$. The effort that maximises the
boats' total profit is $E^{*} = (A - c)/2$. Equilibrium effort exceeds it by the factor $2n/(n + 1)$, and a fee $f =
(A - c)(n - 1)/(2n)$ per boat-day brings the equilibrium to $E^*$.
\end{proposition}

\begin{proof}
Summing the \omterm{def:b2:static-games:best-response}{best responses}, $n(A - E - c) - E = 0$, so $E^{\mathrm{N}} = n(A - c)/(n + 1)$. Total profit $E(A - E - c)$ is
largest at $E^* = (A - c)/2$. The ratio is $2n/(n + 1)$. With a fee $f$, the equilibrium becomes $n(A - c - f)/(n + 1)$;
setting it equal to $(A - c)/2$ gives $A - c - f = (A - c)(n + 1)/(2n)$, so $f = (A - c)(n - 1)/(2n)$.
\end{proof}

Taxing access is one answer. Assigning the resource to an owner, or dividing it into quotas, are others,
and so is leaving the users to manage it themselves. Each needs someone to monitor use and to enforce the
rules, and each has failed somewhere: a fee needs an agency that knows $A$ and $c$ and has reasons to set the
fee right; a quota needs catches to be counted.

\begin{record}[Commons governed by their users]
% ledger: b2:ostrom1990
Ostrom (1990) studied \omterm{def:b2:static-games:common-pool}{common-pool resources}, among them mountain pastures, irrigation systems and
fisheries, that their users had managed over long periods by rules they made themselves, and others where
such rules broke down. The successful cases shared features: clear boundaries of who may use the resource,
rules adapted to local conditions, users who take part in changing the rules, monitors accountable to the
users, sanctions that rise with repeated violations, and cheap ways to settle conflicts. The model of
\cref{prop:b2:static-games:commons} has none of these; its over-use is what happens without them.
\end{record}

\section{Location}\label{sec:b2:static-games:location}

\begin{definition}[Location game]\label{def:b2:static-games:location-game}
In a \emph{location game}\index{location game} two or more sellers (or candidates) choose positions on a line
along which buyers (or voters) are spread; each buyer goes to the nearest seller, and each seller's payoff is
the share of buyers it attracts.
\end{definition}

\begin{omfigure}
\begin{tikzpicture}[econ, font=\footnotesize]
  \draw[thick] (0,0) -- (10,0);
  \foreach \x/\l in {0/0, 10/1} \draw (\x,0.15) -- (\x,-0.15) node[below] {\l};
  \fill[omDef] (3,0) circle (3pt) node[above=4pt] {seller 1 at 0.3};
  \fill[omThm] (6,0) circle (3pt) node[above=4pt] {seller 2 at 0.6};
  \draw[dashed] (4.5,0.5) node[above] {0.45} -- (4.5,-0.5);
  \draw[<->, omDef] (0,-0.9) -- node[below] {share 0.45} (4.5,-0.9);
  \draw[<->, omThm] (4.5,-0.9) -- node[below] {share 0.55} (10,-0.9);
  \draw[->, omDef, thick] (3,1.3) -- (4.3,1.3);
  \node[omDef, anchor=east] at (2.9,1.3) {moves};
\end{tikzpicture}

{\small Two sellers on a line of buyers spread evenly from 0 to 1. Buyers left of the midpoint 0.45 go to seller
1, the rest to seller 2. Seller 1 gains by moving towards seller 2; the only equilibrium has both at the median
buyer, 0.5.}
\end{omfigure}

\begin{proposition}[Two sellers meet at the median]\label{prop:b2:static-games:location}
With buyers spread along a line and two sellers, the unique \omterm{def:b2:static-games:nash-equilibrium}{Nash equilibrium} places both sellers at the median
buyer's position.
\end{proposition}

\begin{proof}
If the sellers are at different positions, each gains by moving towards the other: it keeps every buyer on its
outer side and gains some between them. If both are at the same position away from the median, each has half
the buyers, and either one can move a little towards the median and take more than half. At the median, each
has half; a seller who moves away gets the buyers on its side of the new midpoint, fewer than half.
\end{proof}

With voters instead of buyers, this is the result that two candidates who care only about winning both take
the median voter's position, the companion of the \omterm{def:b2:social-choice:median-voter}{median voter theorem} of \cref{ch:b2:social-choice}. With
three sellers there is no equilibrium in pure strategies (\cref{exo:b2:static-games:12}), and with several
dimensions, as \cref{ch:b2:social-choice} showed, the median may not exist.

\section{Exercises}\label{sec:b2:static-games:exercises}

\begin{exercise}[$\star$]\label{exo:b2:static-games:1}
The row player chooses $T$ or $B$, the column player $L$, $C$ or $R$. Payoffs: $(T, L) = (1, 0)$, $(T, C) = (1, 2)$,
$(T, R) = (0, 1)$, $(B, L) = (0, 3)$, $(B, C) = (0, 1)$, $(B, R) = (2, 0)$. Solve the game by \omterm{def:b2:static-games:iesds}{iterated elimination of
strictly dominated strategies}.
\end{exercise}

\begin{exercise}[$\star$]\label{exo:b2:static-games:2}
Find all pure Nash equilibria of the game with row strategies $U, M, D$ and columns $L, C, R$: row $U$: $(3, 3)$,
$(0, 2)$, $(1, 0)$; row $M$: $(2, 0)$, $(2, 2)$, $(0, 1)$; row $D$: $(0, 1)$, $(1, 0)$, $(4, 4)$.
\end{exercise}

\begin{exercise}[$\star$]\label{exo:b2:static-games:3}
Read the opening chart. By how many points did punishment raise the mean contribution of strangers, and of fixed
groups? Which treatment comes closest to the Nash prediction?
\end{exercise}

\begin{exercise}[$\star$]\label{exo:b2:static-games:4}
Describe the result of a public-goods experiment that would refute the prediction of zero contributions, and
one that would be consistent with it although contributions are positive in early periods.
\end{exercise}

\begin{exercise}[$\star\star$]\label{exo:b2:static-games:5}
With $p = 100 - Q$ and $c = 10$, find each firm's output, the price and each firm's profit for $n = 1, 2, 4, 9$.
\end{exercise}

\begin{exercise}[$\star\star$]\label{exo:b2:static-games:6}
In the public-good game with $y = 12$, find the Nash contributions and the efficient total for $n = 3$.
\end{exercise}

\begin{exercise}[$\star\star$]\label{exo:b2:static-games:7}
Predict, then compute: how does the ratio of the Nash total to the efficient total change from $n = 2$ to $n = 4$ and
to $n = 10$? What does each person give?
\end{exercise}

\begin{exercise}[$\star\star$]\label{exo:b2:static-games:8}
Two boats fish a ground with $A = 100$ and $c = 20$. Find the equilibrium effort, the efficient effort, the
over-use in per cent and the fee that removes it.
\end{exercise}

\begin{exercise}[$\star\star$]\label{exo:b2:static-games:9}
Describe a test of \cref{prop:b2:static-games:location} with data on candidates' positions and voters'
positions in many elections. What result would count against it?
\end{exercise}

\begin{exercise}[$\star\star\star$]\label{exo:b2:static-games:10}
(a) Show that a \omterm{def:b2:static-games:nash-equilibrium}{Nash equilibrium} is never removed by \omterm{def:b2:static-games:iesds}{iterated elimination of strictly dominated strategies}.
(b) Each of $n$ people has utility $x_i + v(G)$, with $v$ increasing and strictly concave, and gives $g_i \ge 0$ from
income $y$. Show that in every \omterm{def:b2:static-games:nash-equilibrium}{Nash equilibrium} with a positive total, $v'(G^{\mathrm{N}}) = 1$, while efficiency requires
$nv'(G^*) = 1$; conclude that $G^{\mathrm{N}} < G^*$ for $n \ge 2$.
\end{exercise}

\begin{exercise}[$\star\star\star$]\label{exo:b2:static-games:11}
Show that matching pennies has no \omterm{def:b2:static-games:nash-equilibrium}{Nash equilibrium} in pure strategies. Then construct a game in which each
player has infinitely many strategies (numbers in $[0, 1]$) and no pure \omterm{def:b2:static-games:nash-equilibrium}{Nash equilibrium} exists.
\end{exercise}

\begin{exercise}[$\star\star\star$]\label{exo:b2:static-games:12}
Three sellers choose positions on $[0, 1]$, with buyers spread evenly; a seller sharing a position with others
splits those buyers equally. Show that there is no \omterm{def:b2:static-games:nash-equilibrium}{Nash equilibrium} in pure strategies.
\end{exercise}

\section{Problem: Kestria's fishing ground}\label{sec:b2:static-games:problem}

\begin{problem}[{Weekend problem --- five boats fish one ground off the coast of Kestria, an invented economy:
how much do they over-fish, and which of two answers, a licence fee or a quota set and monitored by the boats
themselves, restores the efficient catch}]
\label{pb:b2:static-games:1}
A boat-day of effort lands a catch worth $100 - E$ thousand kests, where $E$ is the total effort of all boats in a
season; a boat-day costs 20 thousand kests. There are five boats.

\medskip
\noindent\textbf{Part I --- The catch game.}
\begin{enumerate}
  \item Write boat $i$'s profit as a function of $e_i$ and the others' effort $E_{-i}$.
  \item Derive boat $i$'s \omterm{def:b2:static-games:best-response}{best response}.
  \item Does any boat have a dominant strategy?
  \item Which earlier game of the chapter has the same structure?
  \item Explain in one sentence why a boat ignores part of the cost of its effort.
  \item What would change if the value per boat-day did not depend on $E$?
\end{enumerate}

\medskip
\noindent\textbf{Part II --- The Nash catch.}
\begin{enumerate}[resume]
  \item Find each boat's effort and the total in the symmetric \omterm{def:b2:static-games:nash-equilibrium}{Nash equilibrium}.
  \item Find the value per boat-day and each boat's profit.
  \item Find the boats' total profit.
  \item Check that no boat gains by changing its effort alone.
  \item How would the total effort change with ten boats?
  \item Is the equilibrium \omterm{def:b2:exchange-economies:pareto-efficient}{Pareto-efficient} for the boats? Explain.
\end{enumerate}

\medskip
\noindent\textbf{Part III --- The efficient catch.}
\begin{enumerate}[resume]
  \item Find the total effort that maximises the boats' total profit.
  \item Find the value per boat-day and the total profit there.
  \item By what percentage does the Nash effort exceed it?
  \item How much total profit does over-fishing destroy each season?
  \item The model leaves out the fish stock's future. Would including it make over-fishing worse or better?
  \item Name the result so far.
\end{enumerate}

\medskip
\noindent\textbf{Part IV --- Two answers.}
\begin{enumerate}[resume]
  \item A licence fee of $f$ per boat-day is charged. Find the fee that makes the Nash total equal to the efficient
        effort.
  \item With that fee, find each boat's profit and the fee revenue.
  \item The boats instead agree on a quota of 8 boat-days each. If the other four keep it, what is one boat's best
        effort, and how much would it gain by cheating?
  \item The boats hire a monitor who detects cheating with probability one half. What fine makes the quota hold?
  \item Compare the two answers: what does each need to know, and what can go wrong with each?
  \item State the result: the over-fishing as a share of the efficient effort, and the fee that restores
        efficiency.
\end{enumerate}
\end{problem}
